\documentclass[crop,tikz,11pt]{standalone}
\usepackage{geometricfigures}
\usepackage{amsmath,amssymb}
\usepackage{pgfplots}
\pgfplotsset{compat=1.18}
\usetikzlibrary{arrows.meta,decorations.markings,patterns,intersections,calc}

\definecolor{lib}{HTML}{3B75AF}
\definecolor{rot}{HTML}{C1701E}
\definecolor{sep}{HTML}{B03030}
\definecolor{net}{HTML}{4E9A4E}

% ---------------------------------------------------------------------------
% The cylinder, in the projection and with the conventions of
% pendulum-solution-manifold -- u = theta - pi, so u=0 (the stable equilibrium) is
% at the front and u=+-180 (the unstable one) at the back. Drawn smaller here,
% because it is one panel of three.
\newcommand{\cylR}{1.55}
\newcommand{\cylC}{0.38}
\newcommand{\cylS}{0.78}
\newcommand{\cylP}{2.50}
\pgfmathdeclarefunction{cylx}{1}{\pgfmathparse{\cylR*sin(#1)}}
\pgfmathdeclarefunction{cyly}{2}{\pgfmathparse{-\cylC*\cylR*cos(#1)+\cylS*(#2)}}
\pgfmathdeclarefunction{cylorb}{2}{\pgfmathparse{sqrt(max(0,2*((#2)+cos(#1))))}}

% The canonical chart: theta across, p_theta up, with theta=pi at the centre so
% the two cut edges theta=0 and theta=2pi are the two sides. Angles in degrees.
\newcommand{\canX}{0.86}   % cm per radian of theta
\newcommand{\canY}{0.72}   % cm per unit of p_theta
\pgfmathdeclarefunction{canx}{1}{\pgfmathparse{\canX*((#1)-180)*pi/180}}
\pgfmathdeclarefunction{cany}{1}{\pgfmathparse{\canY*(#1)}}
\pgfmathdeclarefunction{orb}{2}{\pgfmathparse{sqrt(max(0,2*((#2)-cos(#1))))}}

% The three trajectories the picture follows, one of each kind.
% H = 0 turns at theta = pi/2 and 3pi/2, which on the cylinder is exactly the
% silhouette u = +-90: the librating loop is the whole near side and nothing of
% it is hidden. H = 1.4 is the outermost level the paper's figure covers.
\def\Hlib{0.0}
\def\Hrot{1.4}
\def\tlibA{90}
\def\tlibB{270}

% ---------------------------------------------------------------------------
% Panel (ii) is the trained encoder's own output, not a sketch. The curves below
% were produced by pushing analytic level sets of H through the encoder of
%   sae_weights_current.h5
% (GeometricMachineLearning 0.8.0), then simplifying each polyline to 0.008 cm
% with Ramer-Douglas-Peucker. gm/gp are the two separatrix branches, obtained by
% splitting the H -> 1- parametrisation at the unstable equilibrium: s in
% [pi/2,3pi/2] is the p<0 branch l_- and the rest is l_+. Their signed areas are
% -14.7889 and +1.8373 -- opposite orientations, which is why the librating
% limit is the difference, 12.9516. lib is H = 0.0, rot is H = 1.4, the same two
% levels panel (i) and the cylinder draw.
%
% The two axes are scaled differently, 1.5309 cm per unit of z_q against 0.5256
% of z_p, exactly as the paper's own latent panel is: areas in this panel must
% not be compared by eye with panel (i).
% >>> BEGIN GENERATED BLOCK -- generated from the trained network, do not edit by hand
% Weights: sae_weights_current.h5
% Latent extent z_q in [-1.637, 1.172], z_p in [-3.026, 6.677].
% Page scale 1.5309 cm per unit of z_q against 0.5256 of z_p (ratio 2.91): the two axes
% are scaled differently, as the paper's own latent panel is, so areas in this panel
% must not be compared by eye with panel (i).
% Signed areas: gamma_- = -14.7889, gamma_+ = +1.8373, sum = -12.9516.
% lib is H = 0.00, rot is H = 1.40, the same two levels the cylinder and (i) draw.
\def\gmpath{%
  (2.059,-0.738) --
  (2.063,-0.945) --
  (2.014,-1.114) --
  (1.964,-1.197) --
  (1.902,-1.267) --
  (1.740,-1.386) --
  (1.565,-1.459) --
  (1.367,-1.500) --
  (1.182,-1.507) --
  (0.710,-1.475) --
  (0.313,-1.526) --
  (0.182,-1.557) --
  (0.076,-1.609) --
  (-0.036,-1.690) --
  (-0.420,-2.057) --
  (-0.585,-2.192) --
  (-0.763,-2.299) --
  (-0.950,-2.371) --
  (-1.148,-2.412) --
  (-1.360,-2.427) --
  (-1.534,-2.414) --
  (-1.675,-2.373) --
  (-1.786,-2.304) --
  (-1.877,-2.206) --
  (-1.954,-2.071) --
  (-2.014,-1.902) --
  (-2.081,-1.573) --
  (-2.109,-1.207) --
  (-2.090,-0.953) --
  (-2.026,-0.757) --
  (-1.907,-0.595) --
  (-1.719,-0.454) --
  (-1.478,-0.344) --
  (-0.876,-0.142) --
  (-0.736,-0.059) --
  (-0.639,0.043) --
  (-0.581,0.158) --
  (-0.537,0.323) --
  (-0.395,1.446) --
  (-0.353,1.668) --
  (-0.300,1.859) --
  (-0.215,2.065) --
  (-0.111,2.226) --
  (0.010,2.342) --
  (0.144,2.407) --
  (0.290,2.422) --
  (0.450,2.391) --
  (0.633,2.312) --
  (0.840,2.185) --
  (1.116,1.973) --
  (1.403,1.715) --
  (1.630,1.477) --
  (1.776,1.280) --
  (1.840,1.158) --
  (1.890,1.024) --
  (1.963,0.682) --
  (1.993,0.347) --
  (2.000,-0.237) --
  (2.059,-0.738) -- cycle}
\def\gppath{%
  (2.059,-0.738) --
  (2.009,-0.544) --
  (1.972,-0.482) --
  (1.888,-0.407) --
  (1.855,-0.359) --
  (1.760,-0.067) --
  (1.670,0.129) --
  (1.583,0.252) --
  (1.543,0.277) --
  (1.503,0.280) --
  (1.445,0.249) --
  (1.388,0.180) --
  (1.329,0.064) --
  (1.302,-0.034) --
  (1.313,-0.160) --
  (1.411,-0.440) --
  (1.415,-0.499) --
  (1.404,-0.550) --
  (1.338,-0.641) --
  (1.016,-0.938) --
  (0.882,-1.020) --
  (0.728,-1.061) --
  (0.098,-1.075) --
  (-0.200,-1.102) --
  (-0.328,-1.147) --
  (-0.364,-1.181) --
  (-0.372,-1.214) --
  (-0.324,-1.265) --
  (-0.170,-1.335) --
  (-0.016,-1.391) --
  (0.134,-1.417) --
  (0.278,-1.423) --
  (0.550,-1.403) --
  (1.018,-1.425) --
  (1.319,-1.394) --
  (1.481,-1.354) --
  (1.636,-1.296) --
  (1.776,-1.224) --
  (1.890,-1.140) --
  (1.965,-1.060) --
  (2.022,-0.964) --
  (2.052,-0.864) --
  (2.059,-0.738) -- cycle}
\def\libpath{%
  (1.046,-0.528) --
  (1.004,-0.654) --
  (0.941,-0.754) --
  (0.859,-0.825) --
  (0.755,-0.870) --
  (0.569,-0.894) --
  (0.045,-0.876) --
  (-0.234,-0.892) --
  (-0.572,-0.963) --
  (-0.721,-1.013) --
  (-0.822,-1.064) --
  (-0.899,-1.136) --
  (-0.930,-1.223) --
  (-0.924,-1.323) --
  (-0.866,-1.559) --
  (-0.882,-1.675) --
  (-0.932,-1.757) --
  (-1.014,-1.843) --
  (-1.104,-1.910) --
  (-1.199,-1.954) --
  (-1.296,-1.975) --
  (-1.401,-1.972) --
  (-1.508,-1.944) --
  (-1.613,-1.892) --
  (-1.702,-1.825) --
  (-1.779,-1.741) --
  (-1.844,-1.638) --
  (-1.892,-1.522) --
  (-1.946,-1.271) --
  (-1.944,-1.024) --
  (-1.889,-0.834) --
  (-1.837,-0.743) --
  (-1.773,-0.666) --
  (-1.622,-0.543) --
  (-1.431,-0.442) --
  (-1.222,-0.365) --
  (-0.792,-0.240) --
  (-0.636,-0.178) --
  (-0.521,-0.103) --
  (-0.439,-0.006) --
  (-0.392,0.096) --
  (-0.356,0.236) --
  (-0.238,1.092) --
  (-0.152,1.492) --
  (-0.095,1.647) --
  (-0.024,1.773) --
  (0.059,1.863) --
  (0.148,1.913) --
  (0.232,1.927) --
  (0.325,1.911) --
  (0.420,1.867) --
  (0.523,1.791) --
  (0.633,1.681) --
  (0.742,1.544) --
  (0.833,1.402) --
  (0.895,1.268) --
  (0.923,1.174) --
  (0.933,1.081) --
  (0.906,0.751) --
  (0.921,0.339) --
  (0.945,0.179) --
  (1.050,-0.248) --
  (1.060,-0.390) --
  (1.046,-0.524) -- cycle}
\def\rotpath{%
  (2.098,0.597) --
  (2.078,0.999) --
  (2.047,1.138) --
  (1.997,1.262) --
  (1.919,1.389) --
  (1.809,1.526) --
  (1.463,1.861) --
  (1.205,2.074) --
  (0.953,2.257) --
  (0.733,2.394) --
  (0.538,2.488) --
  (0.371,2.538) --
  (0.223,2.549) --
  (0.089,2.521) --
  (-0.031,2.455) --
  (-0.137,2.353) --
  (-0.228,2.220) --
  (-0.308,2.054) --
  (-0.371,1.864) --
  (-0.447,1.491) --
  (-0.585,0.395) --
  (-0.629,0.211) --
  (-0.688,0.082) --
  (-0.783,-0.027) --
  (-0.931,-0.120) --
  (-1.552,-0.339) --
  (-1.795,-0.459) --
  (-1.895,-0.531) --
  (-1.978,-0.613) --
  (-2.043,-0.702) --
  (-2.090,-0.797) --
  (-2.132,-0.946) --
  (-2.149,-1.118) --
  (-2.145,-1.347) --
  (-2.121,-1.617) --
  (-2.082,-1.866) --
  (-2.033,-2.071) --
  (-1.973,-2.232) --
  (-1.898,-2.361) --
  (-1.839,-2.428) --
  (-1.771,-2.481) --
  (-1.697,-2.518) --
  (-1.610,-2.541) --
  (-1.399,-2.545) --
  (-1.110,-2.494) --
  (-0.794,-2.391) --
  (-0.517,-2.243) --
  (-0.286,-2.061) --
  (0.087,-1.691) --
  (0.182,-1.623) --
  (0.274,-1.580) --
  (0.409,-1.549) --
  (0.834,-1.499) --
  (1.376,-1.543) --
  (1.600,-1.540) --
  (1.813,-1.496) --
  (1.982,-1.412) --
  (2.055,-1.346) --
  (2.105,-1.273) --
  (2.135,-1.194) --
  (2.149,-1.097) --
  (2.144,-0.955) --
  (2.088,-0.560) --
  (2.069,-0.313) --
  (2.098,0.589) -- cycle}
\def\pinchX{2.059}\def\pinchY{-0.738}
\def\libtopX{0.232}\def\libtopY{1.927}
\def\gmtopX{0.262}\def\gmtopY{2.423}
\def\gptopX{1.518}\def\gptopY{0.281}
\def\rottopX{0.256}\def\rottopY{2.550}
\def\gplabX{1.050}\def\gplabY{-1.222}
\def\gmleadX{-0.427}\def\gmleadY{1.203}
% <<< END GENERATED BLOCK

\begin{document}
\begin{tikzpicture}[>=Stealth,
  hidden/.style={dash pattern=on 2.2pt off 2.2pt,opacity=0.38},
  map/.style={->,thin,shorten >=2.5pt,shorten <=2.5pt,opacity=0.9},
  panelcap/.style={font=\small,align=center,anchor=north,text width=6.4cm}]

% ===========================================================================
% (0)  M, the solution manifold, on top
% ===========================================================================
\begin{scope}[shift={(0,6.9)}]
  \fill[fg!4!bg]
    plot[domain=-90:90,samples=50,smooth] ({cylx(\x)},{cyly(\x,-\cylP)})
    -- plot[domain=90:270,samples=50,smooth] ({cylx(\x)},{cyly(\x,\cylP)})
    -- cycle;
  \draw[fg!25!bg,densely dotted]
    plot[domain=-180:180,samples=120,smooth] ({cylx(\x)},{cyly(\x,\cylP)});
  \draw[fg!25!bg,densely dotted]
    plot[domain=-180:180,samples=120,smooth] ({cylx(\x)},{cyly(\x,-\cylP)});
  \draw[fg!25!bg] ({-\cylR},{-\cylS*\cylP}) -- ({-\cylR},{\cylS*\cylP});
  \draw[fg!25!bg] ({\cylR},{-\cylS*\cylP}) -- ({\cylR},{\cylS*\cylP});

  % the seam theta=0: the line chart (i) cuts along. It runs up the back.
  \draw[fg!45!bg,dash pattern=on 1.6pt off 1.6pt]
    (0,{\cylC*\cylR-\cylS*\cylP}) -- (0,{\cylC*\cylR+\cylS*\cylP});
  \node[fg!55!bg,font=\scriptsize,anchor=south,inner sep=2pt]
    at (0,{\cylC*\cylR+\cylS*\cylP}) {the seam $\theta=0$, where (i) cuts};

  % --- far side --------------------------------------------------------------
  \draw[rot,thick,hidden]
    plot[domain=90:270,samples=90,smooth] ({cylx(\x)},{cyly(\x,-cylorb(\x,\Hrot))});
  \draw[sep,very thick,hidden]
    plot[domain=90:270,samples=110,smooth] ({cylx(\x)},{cyly(\x,cylorb(\x,1.0))});
  \draw[sep,very thick,hidden]
    plot[domain=90:270,samples=110,smooth] ({cylx(\x)},{cyly(\x,-cylorb(\x,1.0))});

  % --- near side -------------------------------------------------------------
  \fill[lib!16!bg]
    plot[domain=-90:90,samples=60,smooth] ({cylx(\x)},{cyly(\x,cylorb(\x,\Hlib))})
    -- plot[domain=90:-90,samples=60,smooth] ({cylx(\x)},{cyly(\x,-cylorb(\x,\Hlib))})
    -- cycle;
  \draw[lib,thick]
    plot[domain=-90:90,samples=60,smooth] ({cylx(\x)},{cyly(\x,cylorb(\x,\Hlib))});
  \draw[lib,thick]
    plot[domain=-90:90,samples=60,smooth] ({cylx(\x)},{cyly(\x,-cylorb(\x,\Hlib))});
  \draw[sep,very thick]
    plot[domain=-90:90,samples=90,smooth] ({cylx(\x)},{cyly(\x,cylorb(\x,1.0))});
  \draw[sep,very thick]
    plot[domain=-90:90,samples=90,smooth] ({cylx(\x)},{cyly(\x,-cylorb(\x,1.0))});
  \draw[rot,thick]
    plot[domain=-90:90,samples=70,smooth] ({cylx(\x)},{cyly(\x,-cylorb(\x,\Hrot))});

  \fill[sep] (0,{\cylC*\cylR}) circle (1.5pt);

  % where the arrows start: the same three curves, sampled left and right
  \coordinate (Mlib) at ({cylx(-64)},{cyly(-64,-cylorb(-64,\Hlib))});
  \coordinate (Mlm)  at ({cylx(-74)},{cyly(-74,-cylorb(-74,1.0))});
  \coordinate (Mlp)  at ({cylx(-50)},{cyly(-50,cylorb(-50,1.0))});
  \coordinate (Mrot) at ({cylx(-50)},{cyly(-50,-cylorb(-50,\Hrot))});
  \coordinate (Nlib) at ({cylx(64)},{cyly(64,-cylorb(64,\Hlib))});
  \coordinate (Nlm)  at ({cylx(74)},{cyly(74,-cylorb(74,1.0))});
  \coordinate (Nlp)  at ({cylx(50)},{cyly(50,cylorb(50,1.0))});
  \coordinate (Nrot) at ({cylx(50)},{cyly(50,-cylorb(50,\Hrot))});

  \node[lib,font=\scriptsize,anchor=east,inner sep=2pt] at (-1.62,-0.42) {librating};
  \node[sep,font=\scriptsize,anchor=west,inner sep=2pt] at (1.62,1.00) {$\ell_+$};
  \node[sep,font=\scriptsize,anchor=west,inner sep=2pt] at (1.62,-1.34) {$\ell_-$};
  \node[sep,font=\scriptsize,anchor=east,inner sep=2pt] at (-1.62,1.00) {separatrix};
  % Below the rotating orbit, on the cylinder's front, rather than beside it: every point to the
  % right or left of the cylinder at this height is crossed by one of the arrows.
  \node[rot,font=\scriptsize,anchor=north,inner sep=2pt] (rotlab) at (0.55,-2.42) {rotating};
  \draw[rot,thin] (rotlab.north) -- ({cylx(18)},{cyly(18,-cylorb(18,\Hrot))});
  % The panel title sits below the panel, like those of (i) and (ii), in the gap between the arrows.
  \node[font=\small,anchor=north,inner sep=2pt] at (0,-2.98)
    {\textbf{the solution manifold $\mathcal{M}\hookrightarrow\mathbb{R}^4$}};
\end{scope}

% ===========================================================================
% (i)  the canonical chart, bottom left: cut at theta=0, then identify
% ===========================================================================
\begin{scope}[shift={(-4.45,-0.10)}]
  \draw[fg!25!bg] ({canx(0)},-1.78) -- ({canx(360)},-1.78);
  \draw[fg!55!bg,dash pattern=on 1.8pt off 1.8pt] ({canx(0)},-1.78) -- ({canx(0)},1.98);
  \draw[fg!55!bg,dash pattern=on 1.8pt off 1.8pt] ({canx(360)},-1.78) -- ({canx(360)},1.98);
  % the identification, as matching double chevrons on the two cut edges
  \foreach \e in {0,360}{
    \foreach \d in {0,0.14}{
      \draw[fg!55!bg,line width=0.6pt]
        ({canx(\e)-0.11+\d},{-0.80}) -- ({canx(\e)+\d},{-0.66}) -- ({canx(\e)-0.11+\d},{-0.52});
      \draw[fg!55!bg,line width=0.6pt]
        ({canx(\e)-0.11+\d},{0.52}) -- ({canx(\e)+\d},{0.66}) -- ({canx(\e)-0.11+\d},{0.80});
    }
  }
  % That the two dashed edges are one line of M is said in the panel caption: a note above the
  % panel sits where the arrows from M come in.

  % the librating orbit, and the disk it bounds
  \fill[lib!16!bg]
    plot[domain=\tlibA:\tlibB,samples=70,smooth] ({canx(\x)},{cany(orb(\x,\Hlib))})
    -- plot[domain=\tlibB:\tlibA,samples=70,smooth] ({canx(\x)},{cany(-orb(\x,\Hlib))}) -- cycle;
  \draw[lib,thick]
    plot[domain=\tlibA:\tlibB,samples=70,smooth] ({canx(\x)},{cany(orb(\x,\Hlib))})
    -- plot[domain=\tlibB:\tlibA,samples=70,smooth] ({canx(\x)},{cany(-orb(\x,\Hlib))}) -- cycle;
  % the separatrix: both branches, from one cut edge to the other
  \draw[sep,very thick] plot[domain=0:360,samples=120,smooth] ({canx(\x)},{cany(orb(\x,1))});
  \draw[sep,very thick] plot[domain=0:360,samples=120,smooth] ({canx(\x)},{cany(-orb(\x,1))});
  % the rotating orbit: an arc, not a loop
  \draw[rot,thick] plot[domain=0:360,samples=100,smooth] ({canx(\x)},{cany(-orb(\x,\Hrot))});
  \fill[rot] ({canx(0)},{cany(-orb(0,\Hrot))}) circle (1.5pt);
  \fill[rot] ({canx(360)},{cany(-orb(360,\Hrot))}) circle (1.5pt);
  \fill[sep] ({canx(0)},0) circle (1.5pt);
  \fill[sep] ({canx(360)},0) circle (1.5pt);

  \node[font=\scriptsize,anchor=north,inner sep=3pt] at ({canx(0)},-1.78) {$0$};
  \node[font=\scriptsize,anchor=north,inner sep=3pt] at (0,-1.78) {$\pi$};
  \node[font=\scriptsize,anchor=north,inner sep=3pt] at ({canx(360)},-1.78) {$2\pi$};
  \node[font=\scriptsize,anchor=north,inner sep=3pt] at (0,-2.24) {$\theta$};
  \node[font=\scriptsize,rotate=90,anchor=south,inner sep=2pt]
    at ({canx(0)-0.32},0.66) {$p_\theta$};

  \node[sep,font=\scriptsize,inner sep=1.5pt,fill=bg,fill opacity=0.85,text opacity=1]
    at ({canx(215)},1.20) {$\ell_+$};   % right of centre: the arrows arrive from the left
  \node[sep,font=\scriptsize,inner sep=1.5pt,fill=bg,fill opacity=0.85,text opacity=1]
    at ({canx(215)},-1.20) {$\ell_-$};

  % The two angle charts the orbit needs. Chart 1 is the strip itself, theta in (0,2pi); chart 2,
  % theta in (-pi,pi), covers the seam and misses theta = pi instead, so it is drawn as the
  % two halves that the identification glues together. Neither bar covers the whole circle.
  \draw[fg!60!bg,line width=1.1pt] ({canx(0)},-2.80) -- ({canx(360)},-2.80);
  \draw[fg!60!bg,fill=bg] ({canx(0)},-2.80) circle (1.5pt);
  \draw[fg!60!bg,fill=bg] ({canx(360)},-2.80) circle (1.5pt);
  \node[fg!70!bg,font=\scriptsize,anchor=east,inner sep=3pt] at ({canx(0)},-2.80) {chart 1};
  \draw[fg!60!bg,line width=1.1pt,-{Stealth[length=4pt]}] ({canx(180)},-3.08) -- ({canx(360)+0.10},-3.08);
  \draw[fg!60!bg,line width=1.1pt,{Stealth[length=4pt]}-] ({canx(0)-0.10},-3.08) -- ({canx(180)},-3.08);
  \draw[fg!60!bg,fill=bg] ({canx(180)},-3.08) circle (1.5pt);
  \node[fg!70!bg,font=\scriptsize,anchor=east,inner sep=3pt] at ({canx(0)-0.10},-3.08) {chart 2};

  \coordinate (Clib) at ({canx(152)},{cany(orb(152,\Hlib))});
  \coordinate (Cgp)  at ({canx(68)},{cany(orb(68,1))});
  \coordinate (Cgm)  at ({canx(74)},{cany(-orb(74,1))});
  \coordinate (Crot) at ({canx(112)},{cany(-orb(112,\Hrot))});
\end{scope}

% ===========================================================================
% (ii)  the learned chart, bottom right: a plane, with nothing identified
% ===========================================================================
\begin{scope}[shift={(4.45,-0.10)}]
  % what the drawn librating loop bounds, and gamma_+ -- which lies outside it
  \fill[lib!16!bg] \libpath;
  \fill[rot!22!bg] \gppath;

  \draw[rot,thick] \rotpath;
  \draw[lib,thick] \libpath;
  \draw[sep,very thick] \gmpath;
  \draw[sep,very thick] \gppath;
  \fill[sep] (\pinchX,\pinchY) circle (1.6pt);

  \draw[sep!60!bg,line width=0.4pt] (-1.52,1.72) -- (\gmleadX,\gmleadY);
  \node[sep,font=\scriptsize,anchor=east,inner sep=2.5pt] at (-1.52,1.72) {$\gamma_-$};
  \node[sep,font=\scriptsize,inner sep=2pt] at (\gplabX,\gplabY) {$\gamma_+$};
  \node[font=\scriptsize,anchor=west,inner sep=3pt,fg!65!bg] at (2.12,-0.66) {$P$};

  \draw[fg!35!bg,->] (-2.30,-2.62) -- (-1.70,-2.62);
  \draw[fg!35!bg,->] (-2.30,-2.62) -- (-2.30,-2.08);
  \node[font=\scriptsize,anchor=west,inner sep=1.5pt,fg!55!bg] at (-1.68,-2.62) {$z_q$};
  \node[font=\scriptsize,anchor=south,inner sep=1.5pt,fg!55!bg] at (-2.30,-2.06) {$z_p$};

  \coordinate (Llib) at (\libtopX,\libtopY);
  \coordinate (Lgm)  at (\gmtopX,\gmtopY);
  \coordinate (Lgp)  at (\gptopX,\gptopY);
  \coordinate (Lrot) at (\rottopX,\rottopY);
\end{scope}

% ===========================================================================
% the arrows: each trajectory on M to its image in each chart
% ===========================================================================
\draw[map,lib] (Mlib) to[out=205,in=58,looseness=0.8]  (Clib);
\draw[map,sep] (Mlp)  to[out=234,in=52,looseness=0.7]  (Cgp);
\draw[map,sep] (Mlm)  to[out=212,in=96,looseness=0.8]  (Cgm);
\draw[map,rot] (Mrot) to[out=198,in=118,looseness=0.8] (Crot);

\draw[map,lib] (Nlib) to[out=-25,in=122,looseness=0.8] (Llib);
\draw[map,sep] (Nlm)  to[out=-32,in=114,looseness=0.8] (Lgm);
\draw[map,sep] (Nlp)  to[out=-58,in=48,looseness=0.7]  (Lgp);
\draw[map,rot] (Nrot) to[out=-2,in=100,looseness=0.8]  (Lrot);


% The panel titles live in the captions below each panel: anywhere above the panels is crossed by
% the arrows coming down from M.

% ===========================================================================
\node[panelcap] at (-4.45,-3.40)
  {\textbf{(i) the angle chart:} cut along $\theta=0$, then identify the two edges.
   The rotating orbit is an \emph{arc}. It closes only
   because the two edges are the same line of $\mathcal{M}$, so no single angle chart
   holds it: chart 2 covers the seam. $\oint p_\theta\,d\theta=10.13$ is the area
   under it, not an enclosed area.};
\node[panelcap] at (4.45,-3.40)
  {\textbf{(ii) the trained encoder $\Psi^{\mathrm{enc}}$:} one chart, nothing identified, so \emph{every} loop bounds. The
   librating loop bounds the blue region, which leaves $\gamma_+$ out; the rotating
   loop runs outside $\gamma_-$ and takes $\gamma_+$ in. That patch is the step.};

\end{tikzpicture}
\end{document}
